\paragraph{What the ceiling is good for.} A noise ceiling is usually offered as a normalising constant: divide by it, and scores from tasks of different quality become comparable. Our results say it does not work that way on ChEMBL potency data. The fraction of ceiling rises with the very noise it divides out, and it does so because a model evaluated under a random or scaffold split has already seen much of each test molecule's source document. The ceiling is nonetheless useful, as a detector rather than a denominator: a score that approaches or passes it is evidence that the test set carries document support, and the distribution of that support says more about what a number means than the name of the split does. It costs nothing to compute, because ChEMBL already supplies the document identifier.

\paragraph{What this implies for reported split comparisons.} With support held at zero the random, scaffold and temporal splits give the same test error (1.07, 1.07 and 1.06 pChEMBL units). The familiar ordering is therefore, on these data, largely a statement about how many test molecules keep a foothold in their own document, not about generalising to new chemotypes or to the future. We cannot say from one data set how far that carries to the gaps reported elsewhere in the potency-prediction literature, and we do not claim that scaffold splits measure nothing: a scaffold split still removes close analogues, which is a real and different kind of difficulty. What we claim is narrower and checkable --- that a reported split difference on ChEMBL potency data is not interpretable without the support distribution behind it, and that the check is cheap for anyone to run.

\paragraph{What carries the effect.} Of the four mechanisms we could test, three leave no trace (concentration of a task in few documents, between-laboratory offsets inferred from author metadata, potency-dependent heteroscedasticity). The fourth, a constant per-document offset, is real and large --- 57\% of the noise variance --- but bounded: it can be worth at most 0.19 pChEMBL units, against the 0.35 that document support actually delivers. By elimination, the remainder is structural: a document usually reports a series of close analogues, and half a series predicts the other half. We did not run the experiment that would show this directly, which would vary analogue similarity and document membership independently.

\paragraph{What the floor rests on.} Two measurements temper the floor itself. The leave-one-out check shows that the estimate is sound as a second moment but not as a normal law: coverage at the nominal 95\% level is 91.6\%, and tail statements built on $\hat s$ should not be trusted. And the estimate comes from a minority of the material --- two thirds of documents share no molecule with any other, and those documents hold half of all molecules --- which is also the denser and noisier part of the literature. Whether this makes the published floor too high for the isolated half is, strictly, unknowable: the noise of a document with no partner cannot be observed.

\section{Limitations}\label{sec:limits}
{\sloppy
\begin{itemize}
\item \textbf{The floor is cross-document variability, not instrument error.} It includes systematic differences between assays and laboratories and curation errors. Relative to a single well-controlled assay it overstates the irreducible error. Relative to a model that has seen much of a test molecule's own document it is too pessimistic, and Section~\ref{sec:support} quantifies that: the two regimes differ by a factor of two in RMSE. Section~\ref{sec:docsplit} separates them with a document-grouped split.
\item \textbf{Dependence and representativeness of replicates.} The floor is estimated from the 1.8--18.4\% of molecules (median 3.6\%) that were measured in several documents: 8{,}916 molecules out of 210{,}519 over the 41 tasks. These are more potent on average than other molecules in several tasks, are concentrated in few document pairs for some tasks (one pair contributes 79\% of molecule pairs for CHEMBL237 $\mathrm{IC}_{50}$), and include identical values that are probably re-publications. We assume one noise variance per task, and Section~\ref{sec:ceiling} shows that it holds only approximately: $\hat s$ is 0.546 over molecules with exactly two documents against 0.609 over molecules with three or more, higher on 27 of the 41 tasks ($p=0.04$). Section~\ref{sec:provenance} shows that replicates are in fact not independent: document pairs sharing an author agree better than pairs that do not, so the pooled floor is biased downwards. The size of the effect is itself uncertain: it is $\hat\sigma=0.50$ against $0.66$ on the first 17 tasks but only 0.63 against 0.69 on the 24 added ones, and the corrected subsets rest on 1{,}920 of 4{,}992 and 1{,}618 of 7{,}994 pairs respectively, because 40\% and 71\% of pairs involve a document with no author list at all. Assay conditions remain unresolved: ATP concentration is stated in a machine-readable form in only 5--19\% of the kinase assay descriptions, which is too few to correct the Cheng--Prusoff component of the $\mathrm{IC}_{50}$ floors, and we do not substitute literature $K_m$ values in its place. Species and mutant constructs are still not checked.
\item \textbf{The floor is estimated on a minority of the documents, and not a random one.} Only documents that share a molecule with another document contribute to $\hat s$. Over the 41 tasks 8{,}310 of 12{,}339 documents (67.3\%) share none, and the 109{,}597 molecules that appear only in such documents --- 52.1\% of all molecules --- carry labels to which a floor estimated elsewhere is then applied. The estimating minority is also the noisier part: the absolute between-document difference grows with how many partners a document has ($\rho=+0.13$, $p=2\times10^{-46}$ over 12{,}986 differences), and the same direction appears in the dependence on the number of documents per molecule above. Taken literally, both would mean that $\hat s$ is too large for the isolated half of the data and that models are further from the floor than we report, which would make the $f>1$ cases easier rather than harder to explain. We do not claim that: the noise of an isolated document is unobservable in principle, both relations are extrapolations beyond the data, and a simpler reading of each is that a compound gets re-measured because someone is interested in it, in another laboratory and another year. What can be said is that the floor rests on the densest third of the literature.
\item \textbf{Sensitivity to curation choices.} The estimated noise depends on the trimming threshold, on how duplicated values and stereoisomers are treated and on assay selection (Table~\ref{tab:robust}); for some tasks these choices change $\hat s$ by 0.1 or more. The assay confidence threshold of 8 admits homologous proteins.
\item \textbf{Leakage.} Decontamination covers MoleculeNet and TDC but not Polaris datasets or the pretraining corpora of the foundation models, which are large subsets of PubChem and ZINC and may contain test molecules. Such overlap would be expected to favour the pretrained models, which nevertheless perform worse. Removing MoleculeNet and TDC molecules changed test errors very little (Section~\ref{sec:results}).
\item \textbf{What document support does not explain.} The support gradient is measured, not modelled. Of the four candidate mechanisms, three leave no trace in the data (document concentration, between-laboratory offsets, potency-dependent heteroscedasticity) and the fourth, an additive document offset, is now measured and is too small: it accounts for 57\% of the noise variance, which bounds its value at 0.19 pChEMBL units against an observed 0.35 (Sections~\ref{sec:provenance} and~\ref{sec:support}). That leaves close analogues reported together as the remaining explanation, by elimination rather than by measurement. Separating ``learnt the document'' from ``learnt the series'' needs an experiment that varies analogue similarity and document membership independently, which we have not run.
\item \textbf{Model coverage.} We test two small chemical language models, frozen, on all tasks and only a D-MPNN and one fine-tuned model on five tasks with three seeds; larger or more recent foundation models, other heads and longer hyperparameter searches may behave differently. The hyperparameter budget (three configurations) is small for all models.
\item \textbf{Splits.} The temporal split leaks some information (ties at the cut-off year; training labels include later documents), and the inner validation split used for hyperparameters is random for all outer splits --- so hyperparameters are chosen under document support even when the outer split removes it. The document-grouped split assigns each molecule by its earliest document; molecules measured in several documents therefore still share documents across the two sides, and we restrict the support analysis to molecules with a single document for that reason.
\item \textbf{Scope.} Potency endpoints ($K_i$, $\mathrm{IC}_{50}$) from ChEMBL only; we make no claim about ADMET or physicochemical endpoints. $K_i$ and $\mathrm{IC}_{50}$ tasks also differ in target class (mostly GPCRs versus mostly kinases), so differences between them are not attributable to the measurement type alone. Test labels are noisy; we do not correct test metrics for attenuation.
\item \textbf{Literature.} Our account of related work is based on the published abstracts and metadata of several cited papers; we did not reproduce their analyses.
\end{itemize}
\par}

